% !TeX program = tex
\def\ckleapyear#1% check if #1 is a leap year.
{%
  {%
     \def\tleapy{not }% Most years fail.
     \count10=#1
     \count1=#1
     \divide\count1 by 4
     \multiply\count1 by 4
     \ifnum\count1=\count10 % Worrry about 100 & 400.
          \count2=#1
          \divide\count2 by 100
          \multiply\count2 by 100
          \ifnum\count2=\count10 % Not a leap year unless 400 also divides.
               \count3=#1
               \divide\count3 by 400
               \multiply\count3 by 400
               \ifnum\count3=\count10
                    \def\tleapy{}% Special leap year.
               \fi
          \else
               \def\tleapy{}% This gets regular leap years.
          \fi
     \fi
     The year #1 is \tleapy a leap year.\par
  }%
}
\ckleapyear{1900}
\ckleapyear{1990}
\ckleapyear{1996}
\ckleapyear{2000}
\ckleapyear{2004}
\ckleapyear{\the\year}
